% ============================================================
% SECTION 1 — Introduction
% ============================================================
\section{Introduction}
\label{sec:intro}

Flood models produce per-location damage estimates. A crisis manager
uses those estimates for at least six distinct operations: predicting
damage, prioritizing response, identifying where errors concentrate,
transferring models across storms or cities, reasoning about drainage
and adjacency, and allocating resources against forecast lead time.
Each operation can fail for a different spatial reason: prediction can
be inflated by spatial leakage across validation folds; rankings can
shift when infrastructure context enters the evidence; errors can
cluster in neighborhoods the model has not learned; transfer can fail
across events even when within-event accuracy is adequate; adjacency
can be severed by levees or pumps; and temporal evidence can be absent
for events that predate a sensor network.

A single held-out accuracy score addresses prediction and is silent on
the other five. Existing geospatial benchmarks~\cite{yeh2021sustainbench,
klemmer2023satclip, agarwal2025pdfm, krechetova2025geobenchx} report one
metric per task. That metric answers which model ranks best on average;
it does not answer whether the evidence is sufficient for the specific
decision a manager needs to make.

This paper introduces GeoRSCT, a diagnostic benchmark organized around
six spatial \emph{decision geometries}---prediction, ranking, clustering,
transfer, relational propagation, and allocation---each paired with a
named spatial failure mode (\S\ref{sec:problem-geometries}). For each
geometry, the benchmark applies a pre-committed, leakage-controlled test
and reports a sufficiency verdict: \textsc{supported}, \textsc{provisional},
or \textsc{insufficient}. The full set of verdicts is the benchmark's
output: a diagnostic account of what a flood substrate supports and where
the evidence is missing or inadequate.

The benchmark runs on a multi-event flood substrate spanning five U.S.\
metropolitan regions, with radar and satellite precipitation, terrain,
drainage networks, levee and urban infrastructure, insurance claims,
citizen-reported impacts, and high-water marks. Every feature carries
provenance and a field-status annotation so that absent data is
distinguished from a measured zero. The evaluation uses a strict
representation ladder (R0$\to$R1$\to$R2) with fixed folds, solvers, and
seeds; only the feature set changes between levels, so changes in verdict
are attributable to the added evidence. At each level, an RSCT
certificate~\cite{martin2026rsct} decomposes model quality into signal,
suppression, and noise components, providing an auditable record of
\emph{why} a verdict changed, not only that it did. The certificate is
computed and recorded before the next level trains, so the quality
assessment of level $k$ is independent of level $k{+}1$.

GeoRSCT is a benchmark, not a deployed system. It does not propose a
model or optimize a score. Its contribution is a reusable substrate,
a six-geometry evaluation protocol, and a worked example---including
negative results and evidence gaps---of what evidence sufficiency means
in practice for spatial decision support.
